\SourceLeaf{154}
Dass ich die in der Theoria Motus Corporum Coelestium für die Methode der kleinsten Quadrate verwendete Metaphysik später aufgegeben habe, geschah vor allem auch aus einem Grund, den ich selbst öffentlich nicht erwähnt habe. Ich halte es in jeder Hinsicht für weniger wichtig, denjenigen Wert einer unbekannten Größe zu ermitteln, dessen Wahrscheinlichkeit am größten ist, da sie doch immer nur unendlich klein bleibt. Wichtiger ist derjenige Wert, mit dem man sich auf das am wenigsten nachteilige Spiel einlässt. Oder anders gesagt, wenn $fa$ die Wahrscheinlichkeit des Werts $a$ für die Unbekannte $x$ bezeichnet, kommt es weniger darauf an, dass $fa$ ein Maximum wird, als darauf, dass $\int fx \cdot F(x-a)\,dx$, erstreckt über alle möglichen Werte von $x$, ein Minimum wird. Dabei wählt man für $F$ eine Funktion, die immer positiv ist und bei größeren Argumenten in geeigneter Weise immer größer wird. Die Wahl des Quadrats ist rein willkürlich, und diese Willkür liegt in der Natur der Sache. Ohne die bekannten außerordentlich großen Vorteile der Wahl des Quadrats könnte man jede andere Funktion wählen, die diesen Bedingungen genügt. In ganz besonderen Fällen tut man dies auch. Ich weiß aber nicht, ob ich hinreichend klar ausgedrückt habe, was ich mit dieser Unterscheidung eigentlich meine, und bitte andernfalls um freundliche Nachsicht.

. . . . . . . . . . . . . . . .

\Needspace{6\baselineskip}
\begin{center}\textsc{Gauss} an \textsc{Schumacher}. Göttingen, 25. November 1844.\end{center}

. . . . . . . . . . . .

Beim mündlichen Vortrag über die Methode der kleinsten Quadrate pflege ich genau den umgekehrten Weg von dem einzuschlagen, der für eine gedruckte Abhandlung sinnvoll ist. Ich lehre nämlich \textls{zuerst}, wie man sie anwendet, und gehe je nach den Umständen mehr oder weniger auf die feineren Kunstgriffe ein. Erst dann behandle ich, soweit noch Zeit dafür bleibt, die verschiedenen Begründungsarten,
